\hypersetup{pageanchor=false}
\pagenumbering{gobble}
\pagestyle{empty}

\chapter*{Declaration}
\thispagestyle{empty}
I hereby declare that this thesis is my own work and has not been submitted in any form for another degree or qualification at any university or institution.

\clearpage
\chapter*{Acknowledgements}
\thispagestyle{empty}
I would like to thank Dr. Christodoulos Stylianou, my supervisor, for his guidance, constructive feedback, and support during the development of this project. His advice helped me focus the work, improve the technical implementation, and structure the thesis. 

I would also like to thank the Department of Computer Science and Engineering at European University Cyprus for providing the academic environment and resources that supported this work. I am grateful to the lecturers and staff who contributed to my studies and helped me build the knowledge needed to complete this project.

Finally, I would like to thank my family and friends for their encouragement, patience, and support during the preparation of this thesis. Their motivation and understanding were important throughout the research, implementation, and writing process.

\clearpage
\chapter*{Abstract}
\thispagestyle{empty}
\acp{LLM} are becoming increasingly popular in many different \ac{NLP}-related areas (such as question answering, summarization and information retrieval). However, a \ac{LLM} on its own may provide unsupported answers if the required information exists at a remote location or is specific to a particular area of study. \ac{RAG} solves this problem by first finding all relevant parts of a document collection and then making these available to the \ac{LLM} as context prior to generating the final answer. 

This thesis, examines how to optimize and deploy \ac{RAG} pipelines through comparative analysis of six controlled pipeline configurations: a) in memory baseline b) a Milvus vector database pipeline c) a Milvus vector database pipeline with tuning d) dense retrieval with reranking e) hybrid retrieval without reranking f) hybrid retrieval with reranking. For all six of the pipeline configurations described above, we employ the same synthetic enterprise document collection, same set of five scenarios, same log output format and same evaluation procedure. Thus, the only difference among the six pipeline configurations is in their retrieval strategy.

The main comparison uses manual answer score to measure answer correctness and for the retrieval it uses the F1 score to measure how well precision and recall are balanced when the retrieved documents are compared with the expected supporting evidence.

The local results showed that none of the pipelines could have the best F1 scores for both the retrieval and the answers. The pipeline that utilized in-memory data had the highest manual answer mean score, the Milvus based pipeline had the highest average retrieval F1. In some cases reranked search provided better results than a simple first pass, however there was an increase in time to retrieve. Overall these are good indicators that simply having a "better" search engine does not guarantee getting the most accurate answers, as well as all components need to be evaluated to determine if they add sufficient benefits to warrant their costs.

\textcolor{red}{TODO: ADD HPC RESULT HIGHLIGHTS AFTER HPC TESTING.}

The overall framework that has been created in this thesis is a structured local analysis of \ac{RAG} variants using timing-based performance evaluations, quality assessments based on retrieval characteristics, quality assessments of answers generated by models, and a fully reproducible experiment workflow to run and log experiments on both local and \ac{HPC} environment. The full training of models, the fine-tuning of models, the compression of trained models, as well as large scale multi-node benchmarking are beyond the scope of this research.

\clearpage
\chapter*{List of Abbreviations}
\thispagestyle{empty}
\begin{acronym}[FSDP]
\acro{NLP}{Natural Language Processing}
\acro{LLM}{Large Language Model}
\acro{RAG}{Retrieval-Augmented Generation}
\acro{HPC}{High-Performance Computing}
\acro{GPU}{Graphics Processing Unit}
\acro{TPU}{Tensor Processing Unit}
\acro{CPU}{Central Processing Unit}
\acro{LoRA}{Low-Rank Adaptation}
\acro{PEFT}{Parameter-Efficient Fine-Tuning}
\acro{SFT}{Supervised Fine-Tuning}
\acro{ANN}{Approximate Nearest Neighbor}
\acro{KV}{Key--Value}
\acro{FFN}{Feed-Forward Network}
\acro{DPR}{Dense Passage Retrieval}
\acro{BM25}{Best Matching 25}
\acro{HNSW}{Hierarchical Navigable Small World}
\acro{FAISS}{Facebook AI Similarity Search}
\acro{TP}{Tensor Parallelism}
\acro{PP}{Pipeline Parallelism}
\acro{DP}{Data Parallelism}
\acro{FSDP}{Fully Sharded Data Parallel}
\acro{NCCL}{NVIDIA Collective Communications Library}
\acro{MPI}{Message Passing Interface}
\acro{RDMA}{Remote Direct Memory Access}
\acro{TGI}{Text Generation Inference}
\acro{TIS}{Triton Inference Server}
\acro{API}{Application Programming Interface}
\acro{MMLU}{Massive Multitask Language Understanding}
\acro{HELM}{Holistic Evaluation of Language Models}
\acro{MLPerf}{Machine Learning Performance Benchmark Suite}
\acro{PFS}{Parallel File System}
\acro{GPFS}{General Parallel File System}
\acro{IO}{Input/Output}
\acro{TCO}{Total Cost of Ownership}
\acro{OCI}{Open Container Initiative}
\acro{SLURM}{Simple Linux Utility for Resource Management}
\acro{SIF}{Singularity Image Format}
\acro{SFT}{Supervised Fine-Tuning}
\acro{LSTM}{Long Short-Term Memory}
\acro{GRU}{Gated Recurrent Unit}
\acro{GPT}{Generative Pre-trained Transformer}
\acro{RNN}{Recurrent Neural Network}
\acro{RAGAS}{Retrieval Augmented Generation Assessment}
\acro{ROUGE-L}{Recall-Oriented Understudy for Gisting Evaluation - Longest Common Subsequence}
\acro{BERTScore}{Bidirectional Encoder Representations from Transformers Score}
\acro{AWS}{Amazon Web Services}
\acro{I/O}{Input/Output}
\acro{vLLM}{Virtual Large Language Model}
\acro{FP32}{Single-Precision Floating-Point (32-bit)}
\acro{FP16}{Half-Precision Floating-Point (16-bit)}
\acro{INT4}{4-bit Integer Precision}
\acro{INT8}{8-bit Integer Precision}
\acro{GPTQ}{Generalized Post-Training Quantization}
\acro{AWQ}{Activation-aware Weight Quantization}
\acro{KV-cache}{Key-Value Cache}
\acro{KILT}{Knowledge Intensive Language Tasks}
\end{acronym}

\clearpage
\pagenumbering{arabic}
\setcounter{page}{1}
\hypersetup{pageanchor=true}
\pagestyle{frontmatter}

{\setstretch{1.0}
\begingroup

\titleformat{name=\chapter,numberless}
  {\raggedright\bfseries\fontsize{22}{26}\selectfont}
  {}
  {0pt}
  {}

\phantomsection
\markboth{Contents}{}
\tableofcontents

\clearpage
\phantomsection
\addcontentsline{toc}{chapter}{List of Figures}
\markboth{List of Figures}{}
\listoffigures

\clearpage
\phantomsection
\addcontentsline{toc}{chapter}{List of Tables}
\markboth{List of Tables}{}
\listoftables

\endgroup
}

\clearpage
\pagestyle{fancy}
